% theme: chemical_changes_and_batteries
\MajorQuestion{化学変化と電池}
\SetRowSpaceQanda

\Instruction{次の問いに答えなさい。}
\QQ{化学変化によって、物質のもつ化学エネルギーを電気エネルギーに変換する装置を\NamedBlank{battery}という。}{電池（化学電池）}
\QQ{金属が陽イオンになろうとする性質を\NamedBlank{ionization_tendency}という。}{イオン化傾向}
\QQ{うすい塩酸に亜鉛板と銅板を入れて導線でつないだ電池を\NamedBlank{volta_cell}という。}{ボルタ電池}
\QQ{電解質の水溶液に2種類の金属を入れて導線でつなぐと、金属と金属の間に何が生じるか。}{電圧}
\QQ{硫酸亜鉛水溶液と硫酸銅水溶液をセロハンで仕切り、亜鉛板と銅板を入れた電池を\NamedBlank{daniel_cell}という。}{ダニエル電池}

\Instruction{\strut}
\QQ{2種類の金属を用いた\RefBlank{battery}では、\RefBlank{ionization_tendency}が大きい金属が何極になるか。}{負極}
\QQ{\RefBlank{volta_cell}で、電子は亜鉛板と銅板のどちらからどちらへ移動するか。}{亜鉛板から銅板へ移動する}
\QQ{木炭電池で、備長炭とアルミニウムはくの間にはさむ紙にしみこませる液体は何か。}{濃い食塩水}
\QQ{\RefBlank{daniel_cell}で、正極になる金属は何か。}{銅板}
\QQ{亜鉛板を硫酸銅水溶液に入れると、亜鉛板の表面に何が付着するか。}{銅}

\Instruction{\strut}
\QQ{マグネシウム、亜鉛、銅を、\RefBlank{ionization_tendency}が大きい順に並べなさい。}{マグネシウム、亜鉛、銅}
\QQ{\RefBlank{volta_cell}で、負極になる金属は何か。}{亜鉛板}
\QQ{充電してくり返し使うことができる電池を何というか。}{二次電池}
\QQ{\RefBlank{daniel_cell}の正極で起こる変化を、化学式で表しなさい。}{$\mathrm{Cu^{2+} + 2e^- \rightarrow Cu}$}
\QQ{果物を用いた電池で、果物に差し込む2種類の金属板は何か。}{銅板と亜鉛板}

\Instruction{\strut}
\QQ{\RefBlank{volta_cell}の負極で起こる変化を、化学式で表しなさい。}{$\mathrm{Zn \rightarrow Zn^{2+} + 2e^-}$}
\QQ{銅板を硫酸亜鉛水溶液に入れたとき、銅板に変化は見られるか。}{見られない}
\QQ{水の電気分解とは逆の化学変化を利用する電池を何というか。}{燃料電池}
\QQ{\RefBlank{volta_cell}で、銅板の表面から発生する気体は何か。}{水素}
\QQ{2種類の金属を用いた\RefBlank{battery}では、何によって生じる電圧の大きさが変わるか。}{金属の組み合わせ}
